\section{Where liability extends, and where it stops}
\label{sec:16.6}
One of the two bodies of law has now been rewritten for software, and what it covers fits in a sentence: a duty to disclose evidence, and a presumption against the party withholding it, in force where the harm is physical and the defendant is a manufacturer. Those two conditions exclude the cases I have described.

The EU's recast product liability directive (applying from December 2026) makes economic operators liable without fault and carries the instrument built for opacity — a court can order disclosure, and where technical complexity makes proof excessively difficult it may presume defect and causation once the claimant shows each is likely \autocite{eu2024productliability}. A mechanism designed for a system nobody outside can inspect is already law. Two features fix how far it travels: the defendant is that economic operator, not a state deploying what it bought; and the damage is a closed list — death or personal injury including medically recognised psychological harm, property damage, and corruption of non-professional data. The first is wider than it sounds and is where these cases enter if they enter at all: a discriminatory score or an unappealable refusal is recoverable once it has made somebody ill, and not before. The wrong itself, and the economic loss following it, is outside the list.

The proposal that would have covered those was withdrawn — the 2022 AI Liability Directive would have added a fault-based route with the same presumption, beyond defect and beyond the damage list. Member states and industry couldn't agree; it was announced for withdrawal in February 2025 and withdrawn that October \autocite{europeancommission2022liability}. Outside the enacted directive, what remains is Matthias's gap governed by whatever national law says about a case none of it was written for. And a compensation regime of either kind answers one of the four gaps Santoni de Sio and Mecacci distinguish \autocite{santonidesio2021four} — culpability. Moral accountability, public accountability, and the duty to prevent a harm before it occurs are untouched.

The electronic personhood debate approached the gap from the other side, asking whether the system itself could answer. The European Parliament's 2017 resolution floated a status for the most autonomous systems \autocite{europeanparliament2017civil}; in an open letter the following year, 150+ researchers and legal scholars objected that personhood for a system with no consciousness would let a manufacturer offload liability onto a fiction with no assets \autocite{robotics2018openletter}. Giving a bearer assets answers the second half of that objection and not the first. Responsibility follows causal control: the builder answers for how it trained the thing, the operator for what it demanded, the bearer for what it decided knowing what it knew — and no arrangement of accounts moves a design decision away from whoever made it.

Apportioning causal contribution across a training pipeline after a harm has no method, and until it does the apportionment is \emph{drafted}. Two drafts exist: one enacted and confined to defect, one negotiated, withdrawn, and left on the record with its terms intact. Neither was rejected as unworkable, and the second's terms are available to whoever drafts the next one. The nearest proposal aimed at these cases is Crootof's war torts: state liability for civilian harms caused in armed conflict, whether or not anyone committed a war crime \autocite{crootof2016wartorts,crootof2022wartorts}. It covers the state as operator, which the product directive doesn't, and it needs a forum, which the floor court would supply.

Beyond the directive's limits a reviewer can be granted access, publish what it finds, and have produced nothing anybody must act on until somebody is answerable for it.
